%% QAD-MultiGuard — IEEE TMC Submission
%% Revised version: all peer-review issues addressed
%% Template: IEEEtran journal (TMC)
\documentclass[lettersize,journal]{IEEEtran}

% ── Core packages ──────────────────────────────────────────────────────────
\usepackage{amsmath,amsfonts,amssymb}
\usepackage{bm}
\usepackage{algorithmic}
\usepackage{algorithm}
\usepackage{array}
%\usepackage[caption=false,font=normalsize,labelfont=sf,textfont=sf]{subfig}
\usepackage{textcomp}
\usepackage{stfloats}
\usepackage{url}
\usepackage{verbatim}
\usepackage{graphicx}
\usepackage{cite}
\usepackage{booktabs}
\usepackage{multirow}
\usepackage{makecell}
\usepackage[table]{xcolor}
\usepackage{tabularx}
%\usepackage{hyperref}  % conflicts with IEEEtran -- use \url{} from url package

\hyphenation{op-tical net-works semi-conduc-tor IEEE-Xplore quan-ti-za-tion}

% ── Document ────────────────────────────────────────────────────────────────
\begin{document}

\title{QAD-MultiGuard: Multimodal Fraud Detection on Mobile Devices\\
via Quantization-Aware Distillation}

\author{%
  Author information pending%
  \thanks{Author names, affiliations, and funding information to be
    added in the final version.  Manuscript received; revised.}%
}

\markboth{IEEE Transactions on Mobile Computing}%
{Author \MakeLowercase{\textit{et al.}}: QAD-MultiGuard}

\maketitle

% ============================================================================
% ABSTRACT
% ============================================================================
\begin{abstract}
Telecommunication fraud poses a critical challenge to mobile edge
security, where on-device real-time detection must satisfy three
coupled constraints: strict privacy compliance, millisecond-level
latency, and resilience to precision degradation under low-bit
quantization.  We present \textbf{QAD-MultiGuard}, an edge--cloud
collaborative, privacy-preserving multimodal fraud detection
framework that \emph{adapts and extends} the NVFP4 quantization-aware
distillation (QAD) paradigm to the telecommunication anti-fraud
domain.  Our principal algorithmic contribution is the
\textbf{Output-Variance Freeze (OV-Freeze)} regularizer, which
constrains the output-variance drift of quantization-sensitive
$\{q,k,v,o\}$-projection layers to the BF16 teacher baseline via
per-layer scalar rescaling with EMA variance estimation; its gradient
compatibility and joint-loss Lipschitz convergence are formally
analysed in Appendix~A.  Built on top of homologous pure-KL
distillation, OV-Freeze supports a heterogeneous dual-track
quantization scheme---server-side native NVFP4 and on-device
block-wise GGUF-Q4\_K\_M.  For privacy-preserving multimodal fusion,
we construct a 128-dimensional empirically non-reconstructable
acoustic embedding ($\bm{F}_v$, validated against white-box GLO and
black-box model-inversion attacks) and an L-BFGS-optimised
four-modality risk fusion mechanism.  On the TeleAntiFraud-28k
(TAF-28k) benchmark, NVFP4 QAD + OV-Freeze attains $F_1 = 0.923
\pm 0.006$, recovering 99.1\% of the BF16 upper bound and
outperforming ten representative baselines by 6.5--8.5 $F_1$ points.
On the in-house adversarial set AdvFraud-3k the model attains $F_1 =
0.875$, an absolute drop of 0.056 against the BF16 upper bound
(6.0\% relative), well within the deployed threat-model bound
$C_2(\le 6\%)$.  End-to-end deployment on Snapdragon~8~Gen~3 yields
a median per-window alert latency of 268~ms (P99: 342~ms) under
single-call load; under 50 simultaneous calls the P50 degrades to
$\sim$310~ms (P99 $\sim$410~ms), remaining within the 500~ms
incremental budget.
\end{abstract}

\begin{IEEEkeywords}
Telecommunication fraud detection; quantization-aware distillation;
output-variance regularization; privacy preservation; multimodal
fusion; speculative decoding; edge--cloud collaboration.
\end{IEEEkeywords}

% ============================================================================
% 1. INTRODUCTION
% ============================================================================
\section{Introduction}
\IEEEPARstart{T}{he} proliferation of generative AI has driven
telecommunication fraud from script-based single-modality attacks
toward precision-targeted, covert multimodal campaigns.  Ma et
al.~\cite{1} characterised 14 typical fraud-script patterns and 7
acoustic disguise strategies on the TeleAntiFraud-28k corpus, with
each case involving on average 2.7 heterogeneous modality signals.
Single-modality solutions exhibit pronounced detection lag: Wang et
al.~\cite{2} report that ASR-transcript-based detection recall drops
to 60--75\% under dialect, noise, and deliberate voice deformation.

Conventional rule-based or machine-learning approaches cannot
simultaneously satisfy real-time, privacy, and resource constraints.
A typical scam call lasts $\sim$90~s, with critical instructions
concentrated in the first 60~s.  Two latency notions apply: (i)
\emph{per-window incremental latency}---time to complete one
risk-scoring iteration after each sliding window trigger
($\le 500$~ms budget); and (ii) \emph{first-alert latency}---bounded
by the acoustic window length $W$ (first alert within $\sim$3.3~s for
$W=3$~s), after which risk scores refresh at a 1~s cadence.
China's Personal Information Protection Law (PIPL) restricts raw
biometric transfer, mandating on-device processing.  Mobile hardware
imposes strict resource limits (RAM $\le 4$~GB, storage $\le 500$~MB).

Standard post-training quantization (PTQ) applied to 0.5B-scale
models incurs 6.0--12.5\% accuracy degradation~\cite{3}; on TAF-28k
the $F_1$ gap between PTQ and QAD reaches 6.5--8.5 points
(\S\ref{sec:main_results}, Table~\ref{tab3}).  Conventional QAT
requires reconstructing the full SFT$\rightarrow$RL$\rightarrow$merge
pipeline, risking destruction of alignment-stage refusal boundaries.
The NVFP4 QAD report from NVIDIA Nemotron~\cite{3} established that
pure KL-divergence distillation from a BF16 teacher can recover
97.0--99.4\% of full-precision performance, providing the
methodological foundation we build upon.

\subsection{Challenges and Contributions}

Directly applying the native QAD paradigm to anti-fraud detection
faces three challenges:

\textbf{(C1) Limited-data generalisation:} Anti-fraud corpora are
heavily privacy-constrained, requiring empirical verification of
QAD robustness under limited data coverage.

\textbf{(C2) Multimodal fusion under ultra-low-bit quantization:}
Existing QAD work targets single-modality models; hardware
compatibility under NVFP4 and GGUF formats is unexplored.

\textbf{(C3) Edge latency under hardware--software co-design:}
Median alert latency must remain $\le 500$~ms per window; pure
low-bit quantization alone cannot break the mobile compute wall.

This paper makes \textbf{four contributions}:

\textbf{(C1) Adopted methodology---pure-KL homologous distillation
and dual-track quantization.} We \emph{adopt and validate} the
NVFP4 QAD paradigm~\cite{3} on the telecommunication anti-fraud
domain, replacing multi-objective QAT with a single KL loss, and
constructing a heterogeneous dual-track deployment: cloud-side NVFP4
(Blackwell) paired with on-device GGUF-Q4\_K\_M (ARM/Snapdragon).
The domain-adaptation is non-trivial: we show that 0.5B tokens
suffice for full convergence on TAF-28k (Table~\ref{tab6}).

\textbf{(C2) Novel regularizer---Output-Variance Freeze (OV-Freeze).}
We introduce OV-Freeze, a lightweight regularizer that constrains the
output variance of quantization-sensitive self-attention projection
layers ($\{q,k,v,o\}$-proj) to the BF16 teacher baseline via
per-layer scalar rescaling, with EMA variance estimation to suppress
mini-batch noise.  Appendix~A provides formal proofs of gradient
compatibility (Proposition~A.1) and Lipschitz convergence
(Theorem~A.2).  OV-Freeze contributes $+0.7$ $F_1$ points above QAD
alone ($p < 0.01$).

\textbf{(C3) Privacy-preserving multimodal fusion.} We construct a
128-d empirically non-reconstructable acoustic embedding $\bm{F}_v$
validated against GLO and model-inversion attacks ($\text{WER}
\ge 0.95$), and an L-BFGS-optimised four-modality sigmoid fusion
validated via user-stratified 5-fold cross-validation.  Appendix~A.3
provides the non-reconstructability analysis.

\textbf{(C4) Domain-adapted speculative decoding and deployment.}
We fine-tune a 124M-parameter draft model (Qwen2-0.1B) on anti-fraud
corpora, lifting token acceptance rate $\alpha$ from 0.78 to 0.86,
achieving $3.32\times$ speedup on Snapdragon~8~Gen~3.

The remainder is organised as follows. \S\ref{sec:related} reviews
related work; \S\ref{sec:method} details the architecture and
methodology; \S\ref{sec:experiments} reports experiments;
\S\ref{sec:discussion} discusses limitations; \S\ref{sec:conclusion}
concludes.  Appendix~A provides formal analyses.

% ============================================================================
% 2. RELATED WORK
% ============================================================================
\section{Related Work}
\label{sec:related}

\subsection{Quantization-Aware Distillation}

Hinton et al.~\cite{4} introduced knowledge distillation via soft
labels.  Polino et al.~\cite{5} first integrated KD into weight
quantization.  Mishra and Marr~\cite{6} demonstrated that
low-precision networks can match full-precision teachers.  Liu et
al.~\cite{7} validated data-free QAD using model-generated data.  Du
et al.~\cite{8} proposed BitDistiller, combining asymmetric
quantization with self-distillation to mitigate sub-4-bit degradation.

The NVFP4 QAD methodology from NVIDIA Nemotron~\cite{3} is our
primary methodological reference: homologous BF16 self-distillation
outperforms cross-scale configurations (9B$\rightarrow$9B surpasses
12B$\rightarrow$9B); KL-divergence outperforms MSE for distribution
matching; and distillation on math-only corpora simultaneously
recovers coding abilities.  \emph{This paper adopts this QAD
paradigm and extends it with the OV-Freeze regularizer and
multi-modal anti-fraud adaptation.}

Advanced PTQ methods include: AWQ~\cite{9}, protecting critical
channels via activation-aware scaling; GPTQ~\cite{10}, using
second-order Taylor reconstruction; SpinQuant~\cite{11}, suppressing
activation outliers via learnable orthogonal rotations; and
QuaRot~\cite{12}, achieving outlier-free inference via Hadamard
transforms.  These PTQ methods exhibit 6.5--8.5\% accuracy gaps
relative to QAD at 4-bit, which we confirm on TAF-28k
(Table~\ref{tab3}).

\subsection{Telecommunication Fraud Detection}

Early detection~\cite{13,14,15} relied on keyword-matching rules or
shallow ML classifiers.  BERT-Fraud~\cite{16} achieved $F_1 = 0.876$
on a Chinese anti-fraud SMS dataset, but lacks on-device support and
acoustic-modality fusion.

Ma et al.~\cite{1} released TAF-28k, the first open-source
audio-text dual-modality dataset with complex reasoning annotations
(28,511 voice-text pairs, $>$307 hours).  Wang et al.~\cite{2}
proposed SAFE-QAQ, an RL-based audio--text framework (7B parameters,
BF16, server-side inference) achieving $F_1 = 0.918$.  Its reported
production scale of ${\sim}7\!\times\!10^4$ calls/day~\cite{2} is
cited as reported by the authors without independent verification; its
stated principal limitation is the high compute cost precluding
on-device deployment.  \emph{Our work directly addresses this gap:
we achieve comparable $F_1$ (0.923) with a 0.5B quantized model,
achieving ${\sim}28\times$ storage reduction.  The comparison evaluates
a scale--efficiency trade-off rather than a direct algorithmic
superiority claim.}

Table~\ref{tab1} summarises the differentiation.

\begin{table}[htbp]
\caption{Comparison with related work}
\label{tab1}
\centering
\renewcommand{\arraystretch}{1.15}
\small
\setlength{\tabcolsep}{3pt}
\begin{tabular}{>{\raggedright\arraybackslash}p{1.5cm}%
                >{\raggedright\arraybackslash}p{1.75cm}%
                >{\raggedright\arraybackslash}p{1.75cm}%
                >{\raggedright\arraybackslash}p{1.9cm}}
\toprule
\textbf{Dimension} & \textbf{TAF-28k~\cite{1}} &
\textbf{SAFE-QAQ~\cite{2}} & \textbf{This work} \\
\midrule
Core focus     & Dataset + bench  & RL end-to-end    & Edge--cloud sys. \\
Quantization   & N/A              & N/A              & NVFP4 + Q4\_K\_M  \\
Deployment     & ---              & Server (7B)      & Edge ($\le$500 ms) \\
Privacy        & Anonymised       & Unspecified      & $\bm{F}_v$ non-inv. \\
Scale--acc.    & ---              & 7B, $F_1$0.918   & 248 MB, $F_1$0.923 \\
\bottomrule
\end{tabular}
\end{table}

\subsection{On-device LLM Inference}

GGUF~\cite{17} compresses 0.5B models to $\sim$240~MB via Q4\_K\_M
mixed-precision quantization, enabling smooth execution on mobile
ARM SoCs.  vLLM~\cite{18} improves concurrent throughput via
PagedAttention~\cite{19}.  Leviathan et al.~\cite{20} proposed
speculative decoding with theoretical speedup
$(1-\alpha^{\gamma+1})/(1-\alpha)$; Chen et al.~\cite{21} provided
the analytic derivation we follow.

\subsection{Privacy-Preserving Voice Features}

PIPL explicitly mandates that raw biometric data must not leave user
terminals without authorisation.  Tomashenko et al.~\cite{22}
proposed x-vector substitution for acoustic identity obfuscation; Bora
et al.~\cite{23} demonstrated residual invertibility of such embeddings
via GLO attacks.  Our $\bm{F}_v$ construction (dual temporal
compression) aims to sever both semantic and speaker-identity channels
at the physical compression level, as empirically validated in
\S\ref{sec:privacy_exp} and Appendix~A.3.

% ============================================================================
% 3. SYSTEM ARCHITECTURE AND METHOD
% ============================================================================
\section{System Architecture and Method}
\label{sec:method}

\subsection{System Architecture}

\begin{figure}[!t]
\centering
\includegraphics[width=\columnwidth]{Fig/fig1_architecture.png}
\caption{QAD-MultiGuard three-tier edge--cloud collaborative
architecture.  Tier-1 (on-device) extracts the non-invertible
acoustic embedding $\bm{F}_v$ and four scalar risk scores; Tier-2
(cloud) performs NVFP4 chain-of-thought reasoning on guided hidden
states; Tier-3 fuses modality scores via L-BFGS-optimised sigmoid
weights.  Raw audio, SMS plaintext, and ASR transcripts never leave
the device.}
\label{fig1}
\end{figure}

QAD-MultiGuard adopts a three-tier edge--cloud architecture
(Fig.~\ref{fig1}):

\textbf{Tier-1 (On-device detection):} A 240~MB Q4\_K\_M-quantized
student model on Snapdragon~8~Gen~3 constructs a 12-d SMS feature
vector, a 6-d URL structural feature, and the acoustic embedding
$\bm{F}_v \in \mathbb{R}^{128}$.  Raw audio, SMS plaintext, and ASR
transcripts \emph{never leave the device}.  The Whisper-tiny encoder
runs without its decoder head; its pooled output is projected into
the 64-d Whisper branch of $\bm{F}_v$.  The four on-device scalar risk
scores are the only artefacts transmitted to Tier-2.

\textbf{Tier-2 (Cloud deep reasoning):} NVFP4-quantized 0.5B LLM
cluster on Blackwell GPUs performs chain-of-thought (CoT) reasoning
and multimodal cross-validation.  The cloud model operates on guided
hidden states transmitted from the edge---not on raw text---operating
on semantic patterns within the sharply constrained fraud vocabulary
(top-50 terms cover 73.4\% of samples~\cite{1}).  The cloud is
invoked only for medium-confidence cases ($\sim$38\% of calls; \S\ref{sec:exp_setup}).

\textbf{Tier-3 (Multimodal risk fusion):} On-device sigmoid fusion
(\S\ref{sec:fusion}) combining four modality risk scores via
L-BFGS-optimised weights, producing a tri-level risk verdict (Safe /
Medium / High).

\textbf{Edge--cloud coordination.} Calls with on-device score $r
\ge 0.70$ ($\sim$12\%) are intercepted directly; $r \le 0.35$
($\sim$50\%) are cleared on-device; the remaining $\sim$38\% trigger
cloud CoT.  Under 50 simultaneous calls, the cloud batch queue
remains below 150~ms P99 wait time.

\subsection{Formalised Threat Model}

We formalise the threat as $\langle\mathcal{A},\mathcal{K},
\mathcal{G},\mathcal{C}\rangle$:

\textbf{Attacker capability ($\mathcal{A}$):} Active/passive MITM at
the network boundary; can intercept $\bm{F}_v$ and query the cloud
API in black-box fashion; cannot breach the on-device TEE.

\textbf{Attacker knowledge ($\mathcal{K}$):} Full white-box knowledge
of on-device extractors (including $\bm{W}_{\mathrm{proj}}$); black-box
knowledge of cloud topology.

\textbf{Attacker goals ($\mathcal{G}$):} $\mathcal{G}_1$ (privacy
reconstruction from $\bm{F}_v$); $\mathcal{G}_2$ (adversarial evasion
via semantic perturbation); $\mathcal{G}_3$ (identity fraud via
feature replay).

\textbf{Safety constraints ($\mathcal{C}$):}
\begin{itemize}
  \item $C_1$: WER of reconstructed speech $\ge 0.90$.
  \item $C_2$: $F_1$ degradation on AdvFraud-3k $\le 6\%$ relative.
  \item $C_3$: Cross-modal metadata validation thwarts $\mathcal{G}_3$.
\end{itemize}

\subsection{Quantization-Aware Distillation (QAD)}
\label{sec:qad}

\subsubsection{Adopted paradigm: pure-KL homologous distillation}

Following~\cite{3}, we adopt a pure KL-divergence distillation loss:
\begin{equation}
  L_{\mathrm{QAD}} = D_{\mathrm{KL}}\!\bigl(
    p_{\text{teacher}}(y|x)\,\|\,p_{\text{student}}(y|x)\bigr),
  \label{eq:qad}
\end{equation}
with Softmax temperature fixed at $T=1$ throughout distillation.
We use the native Qwen2.5-0.5B-Instruct (BF16)~\cite{27,28,29} as
the homologous teacher---identical architecture to the quantized
student.  As established in~\cite{3}, homologous self-distillation
(9B~BF16 $\rightarrow$ 9B~NVFP4) outperforms cross-scale distillation
(12B~BF16 $\rightarrow$ 9B~NVFP4), because QAD's objective is to
repair output-distribution perturbations from weight discretization,
not to inject new knowledge.  We confirm this on TAF-28k
(Table~\ref{tab6}): our 0.5B homologous teacher converges in 0.5B
tokens, whereas a 7B teacher requires 2.0B tokens for comparable
fidelity.

Conventional QAT retains the SFT cross-entropy, which cannot
constrain non-top-1 token probabilities and causes catastrophic
degradation of RL-stage safety boundaries.  Empirically, QAT yields
KL divergence 0.311 against the BF16 teacher; pure-KL QAD compresses
this to 0.004 (Table~\ref{tab5}).

\subsubsection{Edge--cloud dual-track quantization}

We deploy a heterogeneous dual-track scheme (Table~\ref{tab2}).
Server-side NVFP4~\cite{3} uses block\_size~=~16 with E4M3 FP8 scale
factors; on-device Q4\_K\_M uses 4/6-bit GGUF mixed-precision
optimised for ARM SIMD.  Both tracks are distilled from the same
BF16 teacher.  NVFP4 training and inference are emulated on H100 via
TensorRT-LLM and cross-validated against NVIDIA's public NVFP4
checkpoint (Nemotron-Nano-9B-V2~\cite{3}) with classification-boundary
deviation $\le 0.3\%$.

\begin{table}[htbp]
\caption{Edge--cloud dual-track quantization parameters}
\label{tab2}
\centering
\renewcommand{\arraystretch}{1.15}
\small
\begin{tabular}{l>{\raggedright\arraybackslash}p{2.6cm}%
                  >{\raggedright\arraybackslash}p{2.6cm}}
\toprule
\textbf{Component} & \textbf{Server (NVFP4)} & \textbf{On-device (Q4\_K\_M)} \\
\midrule
Backbone        & Qwen2.5-0.5B-Instruct & Qwen2.5-0.5B-Instruct \\
Parameters      & 494M                  & 494M                   \\
Quantization    & NVFP4 (block=16, E4M3 scale) & Q4\_K\_M (4/6-bit GGUF) \\
Model size      & 248 MB                & 240 MB                 \\
Platform        & Blackwell + vLLM      & ARM / Snapdragon       \\
Teacher         & 0.5B BF16 (homologous)& 0.5B BF16 (homologous) \\
Recovery vs BF16 & \textbf{99.1\%}      & \textbf{98.5\%}        \\
\bottomrule
\end{tabular}
\end{table}

\subsubsection{Output-Variance Freeze (OV-Freeze) regularization}
\label{sec:ovfreeze}

Low-bit discretization causes severe output-variance drift on the
four quantization-sensitive projection layers $\ell \in
\{q,k,v,o\}_{\mathrm{proj}}$, with empirically observed variance
increase up to $+18.2\%$ (Table~\ref{tab:ovf_layer_ablation}).  This
disrupts attention-score scale balance, degrading numerical stability.

OV-Freeze activates only in the \emph{final 30\%} of post-training
optimisation steps ($\ge 600$ of 2,000 steps), after the EMA
variance estimate $\mathrm{Var}^{(t)}_{\mathrm{EMA}}(\bm{y}_\ell)$
has converged near the teacher baseline $\sigma^2_{\mathrm{BF16},\ell}$.
The forward rescaling is:
\begin{equation}
  \bm{y}'_\ell = \bm{y}_\ell \cdot c_\ell,\quad
  c_\ell = \sqrt{\frac{\sigma^2_{\mathrm{BF16},\ell}}
                      {\mathrm{Var}^{(t)}_{\mathrm{EMA}}(\bm{y}_\ell)+\epsilon}},
  \label{eq:ovf}
\end{equation}
where $\epsilon = 10^{-6}$ for numerical stability.  The EMA update is:
\begin{equation}
  \mathrm{Var}^{(t)}_{\mathrm{EMA}} = \rho\,\mathrm{Var}^{(t-1)}_{\mathrm{EMA}}
    + (1-\rho)\,\mathrm{Var}^{(t)}_{\mathrm{batch}},\quad \rho = 0.95.
  \label{eq:ema}
\end{equation}
The EMA buffer is initialised to $\sigma^2_{\mathrm{BF16},\ell}$.  With
$\rho = 0.95$ the effective half-life is $\sim$14 steps ($\sim$110
independent samples), suppressing variance-estimate CV below 0.05.

\textbf{Gradient compatibility.} In the forward--backward graph the
variance branch is detached, making $c_\ell$ a constant scaling
factor:
\begin{equation}
  \frac{\partial L_{\mathrm{QAD}}}{\partial\bm{y}_\ell}
    = c_\ell \cdot \frac{\partial L_{\mathrm{QAD}}}{\partial\bm{y}'_\ell}.
  \label{eq:ovf-detach}
\end{equation}
This preserves the KL-gradient direction exactly, introducing only
per-layer magnitude scaling.  Appendix~A proves that $c_\ell \in
[c_{\min},c_{\max}]$ is bounded away from zero and infinity
(Proposition~A.1) and that the joint loss retains Lipschitz-continuous
gradient convergence (Theorem~A.2).

\textbf{Relationship to layer normalisation.} OV-Freeze differs
from BatchNorm/LayerNorm in two key ways: (i) it targets the
\emph{output} of existing projection layers without modifying the
attention topology or adding new structural parameters; and (ii) the
target variance $\sigma^2_{\mathrm{BF16},\ell}$ is fixed offline from
the BF16 teacher calibration, not learned from training data.
Inserting LayerNorm after each projection layer would introduce
additional parameters and would interact with quantization noise
differently; ablation experiments in
Table~\ref{tab:ovf_layer_ablation} confirm that OV-Freeze on the
four attention projections produces the largest PPL reduction
($\Delta\text{PPL} = -0.11$) at zero parameter overhead.

\subsection{Privacy-Preserving Acoustic Embedding}
\label{sec:acoustic}

The acoustic embedding is:
\begin{equation}
  \bm{F}_v = \mathcal{C}\!\left(\bm{f}_{\mathrm{mfcc}},\,
    \psi(\bm{W}_{\mathrm{proj}}\bar{\bm{h}}_w)\right) \in \mathbb{R}^{128},
  \label{eq:fv}
\end{equation}
where $\mathcal{C}$ denotes concatenation and $\psi$ row-wise
normalisation.  The MFCC vector $\bm{f}_{\mathrm{mfcc}} \in
\mathbb{R}^{64}$ is extracted at 16~kHz, 64-band mel filter banks,
10~ms hop, 25~ms window.  The pooled Whisper-tiny encoder
representation $\bar{\bm{h}}_w \in \mathbb{R}^{384}$ is projected to
64-d via learnable $\bm{W}_{\mathrm{proj}} \in \mathbb{R}^{64\times384}$,
jointly optimised in the QAD phase with row-wise normalisation
preventing projection collapse.

\textbf{Non-reconstructability: engineering analysis.}  We emphasise
that the privacy guarantee is \emph{empirical} rather than
information-theoretic: because $\bm{F}_v$ retains task-relevant
acoustic statistics (energy variance, pitch range), a small but
non-zero mutual information $I(x;\bm{F}_v)>0$ necessarily exists.
The practical protection arises from dual temporal compression: (i)
$\sim$300 MFCC frames are mean-pooled into a single 64-d vector,
discarding temporal ordering and phase; (ii) Whisper global
average-pooling combined with a rank-limited projection
($\mathrm{rank}(\bm{W}_{\mathrm{proj}})\le 64$) severs token-level
fine-grained context.  Empirically, both white-box GLO and black-box
model-inversion attacks yield $\text{WER}\ge 0.95$ and speaker-ID
accuracy $\le 8.3\%$ (approaching the 10\% random baseline),
indicating the residual $I(x;\bm{F}_v)$ is insufficient for any known
inverse attack to recover intelligible speech.  Full analysis in
Appendix~A.3.

\subsubsection{Acoustic proxy statistics for $r_{\mathrm{audio}}$}

Four proxy statistics are derived from $\bm{F}_v$ to form the
acoustic risk score.  Let $e_t = \sum_k f_{\mathrm{mel}}(t,k)^2$
denote per-frame mel energy and $m_t = \overline{\bm{f}_{\mathrm{mel}}(t,0\!:\!8)}$
the low-band mean:
\begin{align}
  \texttt{energy\_var}    &= \mathrm{Var}\bigl(\{e_t\}_{t=1}^{T}\bigr),
    \label{eq:ev} \\
  \texttt{tone\_proxy}    &= \mathbb{E}\bigl[\,\bigl|\bm{f}_{\mathrm{mel}}[\cdot,16\!:\!32]\bigr|\,\bigr],
    \label{eq:tp} \\
  \texttt{urgency\_proxy} &= \mathbb{E}\bigl[\,|\Delta e_t|\,\bigr],
    \label{eq:up} \\
  \texttt{pitch\_range}   &= \max_t m_t - \min_t m_t.
    \label{eq:pr}
\end{align}
Here \texttt{energy\_var} captures speech-rate anomalies;
\texttt{tone\_proxy} is the mid-band ($16$--$32$) magnitude mean,
capturing elevated mid-frequency energy indicative of
official-impersonation scripts; \texttt{urgency\_proxy} is the mean
energy-gradient magnitude, capturing abrupt energy changes correlated
with coercion language; and \texttt{pitch\_range} captures low-band
prosodic deviations.  These four features are linearly aggregated as
$r_{\mathrm{audio}} = \sum_{i=1}^4 \alpha_i s_i$ with
$\bm{\alpha}=[0.35,0.28,0.25,0.12]^T$ learned via $L_2$-regularised
logistic regression ($\lambda=0.01$) on Z-score-standardised features.

\subsubsection{Streaming sliding-window implementation}
\label{sec:streaming}

A $W=3$~s ring buffer advances in $\Delta=1$~s steps (66.7\% overlap).
Per-window P50/P99 latency decomposes as:
(a)~on-device extraction ($\sim$12~ms);
(b)~encrypted upload ($\sim$5~ms RTT);
(c)~cloud NVFP4 CoT with speculative decoding ($\sim$230~ms);
(d)~score download and fusion ($\sim$21~ms).
Total P50 = 268~ms under single-call load on a lab 5G/Wi-Fi~6 network.
Under 50 simultaneous calls, P50 degrades to $\sim$310~ms (P99
$\sim$410~ms), remaining within the 500~ms incremental budget.  The
first-alert latency is $W + t_{\mathrm{infer}} \approx 3.3$~s, after
which scores refresh at a 1~s cadence.

\subsection{Domain-Adapted Speculative Decoding}

Anti-fraud text has high pattern-density; the top-50 fraud terms
cover 73.4\% of TAF-28k~\cite{1}.  We fine-tune a 124M-parameter draft
model (Qwen2-0.1B) on 50K anti-fraud utterances to lift token
acceptance rate $\alpha$ from 0.78 to 0.86.  Expected yield per
iteration: $(1-\alpha^{\gamma+1})/(1-\alpha)$.  At $\gamma=5$:
generic $\alpha=0.78 \Rightarrow 3.52$ tokens; domain-tuned
$\alpha=0.86 \Rightarrow 4.25$ tokens.  The theoretical $4.25\times$
speedup is offset by KV-cache and verification overheads
($\eta_{\text{H100}}\approx 0.82$, $\eta_{\text{SD8G3}}\approx 0.78$),
yielding measured $3.49\times$ (H100) and $3.32\times$ (SD8G3).
$\gamma=5$ is the Pareto optimum on memory-constrained mobile hardware
(details in \S\ref{sec:speculative_exp}).

\subsection{Multimodal Risk Fusion}
\label{sec:fusion}

The global risk score is:
\begin{equation}
  r = \sigma\!\left(\sum_{m} w_m r_m + b\right),
  \label{eq:fusion}
\end{equation}
where $m \in \{\text{text, audio, url, meta}\}$ and $\sigma$ is
sigmoid.  Parameters $\mathbf{w}$ and $b$ are optimised via L-BFGS on
a labelled validation subset.  \emph{User-stratified} 5-fold
cross-validation (all calls from one user in the same fold) yields
$\mathbf{w} = [0.40\pm0.018, 0.30\pm0.015, 0.20\pm0.012,
0.10\pm0.009]^T$.  The intra-modality weights $\bm{\alpha}$ for
$r_{\mathrm{audio}}$ (Eq.~\eqref{eq:ev}--\eqref{eq:pr}) are learned
separately via $L_2$-regularised logistic regression, decoupling
interpretability from cross-modal global optimality.

Risk thresholds: $r\ge 0.70 \Rightarrow$ High (intercept);
$0.35\le r < 0.70 \Rightarrow$ Medium; otherwise Safe.  High-risk
interception requires two consecutive scores $r \ge 0.70$
(inter-window EMA: $r_t = 0.4\,r_{t-1} + 0.6\,r_t^{(\mathrm{new})}$).

\textbf{Fusion strategy comparison.} A 2-layer multimodal Transformer
baseline (cross-attention over four modality tokens) yields only
$+0.004$ $F_1$ gain ($0.927$ vs.\ $0.923$) at $42\times$ parameter
inflation and $+7.2$~ms latency.  Under edge constraints, linear
sigmoid fusion approaches the deep-fusion ceiling at negligible cost.

\textbf{Cross-domain weight stability.} The fusion weights were
optimised on the TAF-28k validation subset and applied unchanged to
ChiFraud and AdvFraud-3k evaluation.  The low cross-fold variance
($\sigma_w \le 0.018$) indicates the learned weights generalise across
the same domain; cross-domain recalibration is left to future work.

% ============================================================================
% 4. EXPERIMENTAL EVALUATION
% ============================================================================
\section{Experimental Evaluation}
\label{sec:experiments}

\subsection{Experimental Setup}
\label{sec:exp_setup}

\subsubsection{Datasets}

\textbf{TAF-28k~\cite{1}:} 28,511 voice-text pairs ($\sim$307 h),
split 8:1:1.  Publicly available on HuggingFace
(\texttt{JimmyMa99/TeleAntiFraud}).  Evaluation follows the
TeleAntiFraud-Bench protocol.

\textbf{AdvFraud-3k:} 1,000 TAF-28k fraud-text samples manually
rewritten under 8 adversarial strategies (synonym substitution,
syntactic reordering, dialect shift, metaphor, topic hijacking,
emotional manipulation, authority escalation, urgency fabrication)
plus 2,000 expert-composed novel fraud scripts.  Cohen's $\kappa =
0.87$ (three annotators, one adjudicator).  Used exclusively for
evaluation; no train-test leakage.  The 1,000 rewritten samples share
vocabulary distribution with training data; they test robustness to
surface perturbations, while the 2,000 novel scripts test
out-of-template generalisation.  Construction protocol available in
supplementary material.

\textbf{ChiFraud~\cite{chifraud}:} 411,934 Chinese network-fraud
texts (59,106 fraud + 352,828 normal), collected 2022--2023; used as
an out-of-distribution (OOD) cross-domain generalisation benchmark.
Publicly available (\texttt{xuemingxxx/ChiFraud}).

\subsubsection{Baselines}

PTQ methods (all NVFP4, without distillation): AWQ, GPTQ, SpinQuant,
QuaRot, and an additional ``NVFP4 PTQ + temperature scaling'' control
(Softmax temperature calibrated on 100 held-out samples to minimise
KL divergence against the BF16 teacher).  Self-distillation:
BitDistiller.  QAT: NVFP4 cross-entropy QAT.  Edge PTQ: Q4\_K\_M
PTQ.  Domain-specific: BERT-Fraud, SAFE-QAQ.

\textbf{Note on PTQ vs.\ QAD comparison.}  PTQ methods are
post-training techniques not designed for distribution alignment; the
observed gap measures the benefit of \emph{distillation-aware}
optimisation.  Temperature scaling as a post-hoc control improves PTQ
$F_1$ from 0.838 to 0.852 but remains 7.1 points below QAD+OVF,
confirming that accuracy recovery derives from distillation-driven
distribution alignment rather than simple temperature rescaling.

\textbf{Note on SAFE-QAQ.} SAFE-QAQ (7B params, BF16, server-side)
is not quantized for this comparison because its RL-based multi-stage
training pipeline is not designed for low-bit deployment; adapting it
would require substantial architectural modification beyond the scope
of this paper.  The comparison evaluates a scale--efficiency
trade-off: our 248~MB model vs.\ a 7B unquantized production system.
A fair same-scale comparison (0.5B SAFE-QAQ under identical
quantization) is left to future work.

\subsubsection{Hardware Platforms}

\textbf{Server:} NVIDIA H100 80GB SXM5; NVFP4 emulated via
TensorRT-LLM.

\textbf{Edge:} Snapdragon~8~Gen~3 (Hexagon NPU, Cortex-X4, 12~GB RAM).

\textbf{Training:} $8\times$H100, PyTorch~2.5, DeepSpeed ZeRO-3.

\subsubsection{Training Configuration}

Cosine LR decay, initial $1\times10^{-5}$, 100 warm-up steps, batch
size 8, context 4096, 2,000 steps, $\sim$0.5B tokens.  OV-Freeze
activates at step 1,400 (final 30\%) with $\lambda=1.0$.  All
experiments: 5 random seeds; mean $\pm$ std reported.  Key pairwise
comparisons: two-sided paired bootstrap ($10^4$ resamples, $p<0.01$).

\subsection{Main Results}
\label{sec:main_results}

\begin{table*}[!t]
\caption{TAF-28k test-set multimodal fraud detection results.
$^\dagger$Temperature scaling: Softmax temperature tuned on 100
held-out samples to minimise KL divergence against the BF16 teacher.}
\label{tab3}
\centering
\renewcommand{\arraystretch}{1.2}
\small
\begin{tabularx}{\textwidth}{Xccccc}
\toprule
\textbf{Method} & \textbf{$F_1$} & \textbf{Precision} & \textbf{Recall}
  & \textbf{FPR} & \textbf{Recovery (\%)} \\
\midrule
BF16 (upper bound) & $0.931\pm0.005$ & $0.928$ & $0.934$ & 1.6\% & 100.0 \\
\midrule
\multicolumn{6}{l}{\emph{Advanced PTQ (all NVFP4 format, without distillation)}} \\
NVFP4 PTQ          & $0.838\pm0.011$ & $0.847$ & $0.829$ & 2.8\% & 90.0 \\
NVFP4 + AWQ~\cite{9}        & $0.838\pm0.010$ & $0.846$ & $0.830$ & 2.7\% & 90.0 \\
NVFP4 + GPTQ~\cite{10}      & $0.840\pm0.010$ & $0.848$ & $0.832$ & 2.7\% & 90.2 \\
NVFP4 + SpinQuant~\cite{11} & $0.838\pm0.011$ & $0.846$ & $0.830$ & 2.7\% & 90.0 \\
NVFP4 + QuaRot~\cite{12}    & $0.838\pm0.011$ & $0.846$ & $0.830$ & 2.7\% & 90.0 \\
NVFP4 PTQ + temp.\ scaling$^\dagger$ & $0.852\pm0.010$ & $0.858$ & $0.846$ & 2.5\% & 91.5 \\
\midrule
\multicolumn{6}{l}{\emph{Self-distillation}} \\
NVFP4 + BitDistiller~\cite{8} & $0.858\pm0.009$ & $0.866$ & $0.850$ & 2.5\% & 92.2 \\
\midrule
\multicolumn{6}{l}{\emph{QAT / QAD (ours --- adapted from~\cite{3})}} \\
NVFP4 QAT          & $0.844\pm0.014$ & $0.853$ & $0.835$ & 3.1\% & 90.7 \\
NVFP4 QAD (no OVF) & $0.916\pm0.007$ & $0.918$ & $0.914$ & 1.9\% & 98.4 \\
\rowcolor[gray]{0.93}\textbf{NVFP4 QAD + OV-Freeze (ours)}
  & $\mathbf{0.923\pm0.006}$ & $\mathbf{0.925}$ & $\mathbf{0.921}$
  & $\mathbf{1.8\%}$ & $\mathbf{99.1}$ \\
\midrule
\multicolumn{6}{l}{\emph{Edge Q4\_K\_M quantization (ours)}} \\
Q4\_K\_M PTQ              & $0.851\pm0.011$ & $0.864$ & $0.838$ & 2.6\% & 91.4 \\
Q4\_K\_M QAD (no OVF)     & $0.911\pm0.008$ & $0.913$ & $0.909$ & 2.0\% & 97.9 \\
\rowcolor[gray]{0.93}\textbf{Q4\_K\_M QAD + OV-Freeze (ours)}
  & $\mathbf{0.917\pm0.007}$ & $\mathbf{0.920}$ & $\mathbf{0.914}$
  & $\mathbf{1.9\%}$ & $\mathbf{98.5}$ \\
\midrule
\multicolumn{6}{l}{\emph{Domain-specific baselines}} \\
BERT-Fraud~\cite{16}     & $0.876$ & $0.882$ & $0.870$ & 2.3\% & --- \\
SAFE-QAQ~\cite{2} (7B, BF16, unquantized)
  & $0.918\pm0.006$ & $0.921$ & $0.916$ & 1.8\% & --- \\
\bottomrule
\end{tabularx}
\end{table*}

Table~\ref{tab3} shows four key findings.  \textbf{First}, NVFP4 QAD
+ OV-Freeze attains $F_1=0.923\pm0.006$ (recovery 99.1\%), outperforming
PTQ methods by $0.065$--$0.085$ absolute $F_1$ points and BitDistiller
by $0.065$.  Temperature scaling as a post-hoc control reaches only
$F_1=0.852$, confirming that distillation-driven distribution
alignment is the primary recovery mechanism.  \textbf{Second}, QAD
($F_1=0.916$, recovery 98.4\%) outperforms QAT ($F_1=0.844$, recovery
90.7\%) by 7.2 points.  \textbf{Third}, OV-Freeze contributes
$+0.7$ recovery points and reduces FPR to 1.8\% ($p<0.01$, bootstrap
95\% CI [$+0.31$, $+1.08$]).  \textbf{Fourth}, the edge Q4\_K\_M
track trails cloud NVFP4 by only $0.006$ $F_1$, confirming
cross-format portability.

Against SAFE-QAQ ($F_1=0.918$, 7B params, BF16, unquantized):
our model achieves comparable accuracy with $\sim$28$\times$ storage
reduction (248~MB vs.\ $\sim$7~GB).  This is a \emph{scale--efficiency
trade-off} comparison, not an algorithmic superiority claim; a
fair same-scale comparison with a quantized SAFE-QAQ remains future
work.

\subsection{Cross-dataset Generalisation and Adversarial Robustness}

\begin{table}[htbp]
\caption{Cross-dataset and adversarial robustness results}
\label{tab4}
\centering
\renewcommand{\arraystretch}{1.15}
\small
\setlength{\tabcolsep}{3pt}
\begin{tabular}{>{\raggedright\arraybackslash}p{0.85cm}%
                >{\raggedright\arraybackslash}p{2.25cm}%
                >{\centering\arraybackslash}p{0.65cm}%
                >{\centering\arraybackslash}p{0.65cm}%
                >{\centering\arraybackslash}p{1.15cm}%
                >{\centering\arraybackslash}p{0.85cm}}
\toprule
\textbf{Train} & \textbf{Test} & \textbf{BF16}
  & \textbf{PTQ} & \textbf{QAD+OVF} & \textbf{$\Delta F_1$} \\
\midrule
TAF-28k & TAF-28k (IID)         & 0.931 & 0.838 & \textbf{0.923} & $-0.008$ \\
TAF-28k & AdvFraud-3k (adv.)   & 0.882 & 0.778 & \textbf{0.875} & $-0.007$ \\
TAF-28k & ChiFraud (OOD)       & 0.871 & 0.768 & \textbf{0.860} & $-0.011$ \\
TAF-28k+CF & AdvFraud-3k (multi) & 0.889 & 0.789 & \textbf{0.881} & $-0.008$ \\
\bottomrule
\end{tabular}
\end{table}

AdvFraud-3k drop quantification: absolute drop vs.\ BF16
AdvFraud-3k baseline ($F_1=0.882$) is $0.007$; relative drop vs.\
TAF-28k BF16 upper bound ($F_1=0.931$) is $6.0\%$ (absolute $0.056$);
relative drop vs.\ deployed model ($F_1=0.923$) is $5.2\%$ (absolute
$0.048$).  Both relative figures satisfy threat-model constraint
$C_2(\le6\%)$.  The 1,000 rewritten AdvFraud samples test surface
robustness; the 2,000 novel scripts test out-of-template
generalisation.  Joint TAF-28k + ChiFraud training lifts adversarial
$F_1$ to $0.881$.

\subsection{Loss Function Ablation}

\begin{table}[htbp]
\caption{Ablation on QAD alignment loss}
\label{tab5}
\centering
\renewcommand{\arraystretch}{1.25}
\small
\setlength{\tabcolsep}{4pt}
\begin{tabular}{lccc}
\toprule
\textbf{Loss variant} & \textbf{$F_1$} & \textbf{KL vs.\ BF16} & \textbf{Rec.\ (\%)} \\
\midrule
\textbf{Pure KL (from~\cite{3})} & $\mathbf{0.916\pm0.007}$ & $\mathbf{0.005}$ & $\mathbf{98.4}$ \\
Logits MSE           & $0.901\pm0.010$ & $0.082$ & $96.8$ \\
Cross-entropy (QAT)  & $0.844\pm0.014$ & $0.311$ & $90.7$ \\
Three-term mixed     & $0.879\pm0.012$ & $0.124$ & $94.4$ \\
KL + task reg.       & $0.908\pm0.009$ & $0.041$ & $97.5$ \\
\bottomrule
\end{tabular}
\end{table}

The ``KL + task regularization'' variant appends a cross-entropy term
for the final 600 steps.  The counterintuitive degradation
($0.916\rightarrow0.908$) arises because hard CE labels pull
non-top-1 probabilities toward a degenerate one-hot target, increasing
KL divergence against the teacher ($0.005\rightarrow0.041$) and
undoing distributional fidelity built during pure-KL pre-training.
This reinforces the finding in~\cite{3} that any hard-label auxiliary
loss creates a gradient conflict with the KL objective on non-maximum
tokens under low-bit quantization constraints.

\subsection{Teacher Selection Ablation}

\begin{table}[htbp]
\caption{Teacher knowledge-transfer ablation}
\label{tab6}
\centering
\renewcommand{\arraystretch}{1.15}
\small
\setlength{\tabcolsep}{4pt}
\begin{tabular}{l>{\centering\arraybackslash}p{1.05cm}%
                 >{\centering\arraybackslash}p{1.0cm}%
                 >{\centering\arraybackslash}p{1.05cm}%
                 >{\centering\arraybackslash}p{1.0cm}}
\toprule
\textbf{Teacher} & \textbf{$F_1$ @0.5B tok} & \textbf{$F_1$ conv.} & \textbf{Tok.\ to conv.} & \textbf{Rec.\ (\%)} \\
\midrule
\textbf{0.5B (homol.)} & $\mathbf{0.916}$ & $\mathbf{0.916}$ & $\mathbf{0.5\text{B}}$ & $\mathbf{98.4}$ \\
1.8B BF16 & $0.911$ & $0.913$ & $0.7\text{B}$ & $97.9$ \\
3B BF16   & $0.904$ & $0.910$ & $1.0\text{B}$ & $97.1$ \\
7B BF16   & $0.892$ & $0.915$ & $2.0\text{B}$ & $95.8$ \\
\bottomrule
\end{tabular}
\end{table}

Homologous 0.5B teacher achieves peak $F_1$ at convergence with
$4\times$ less data (0.5B vs.\ 2.0B tokens) than a 7B teacher,
consistent with~\cite{3}.

\subsection{OV-Freeze Ablation}
\label{sec:ovf_ablation}

\begin{table}[htbp]
\caption{OV-Freeze layer-placement ablation}
\label{tab:ovf_layer_ablation}
\centering
\renewcommand{\arraystretch}{1.15}
\small
\begin{tabular}{>{\raggedright\arraybackslash}p{3.2cm}cccc}
\toprule
\textbf{Strategy} & \textbf{$F_1$} & \textbf{PPL} & \textbf{$\Delta$PPL} & \textbf{Var.\ drift} \\
\midrule
No OVF baseline      & $0.916$ & $8.73$ & ---    & $+18.2\%$ \\
+ OVF (FFN only)     & $0.918$ & $8.71$ & $-0.02$ & $+15.4\%$ \\
+ OVF ($q$-proj)     & $0.918$ & $8.69$ & $-0.04$ & $+9.4\%$  \\
+ OVF ($q,v$)        & $0.920$ & $8.66$ & $-0.07$ & $+5.1\%$  \\
+ OVF ($q,k,v$)      & $0.922$ & $8.64$ & $-0.09$ & $+2.8\%$  \\
\textbf{+ OVF ($q,k,v,o$)} & $\mathbf{0.923}$ & $\mathbf{8.62}$
  & $\mathbf{-0.11}$ & $\mathbf{+1.3\%}$ \\
+ OVF ($q,k,v,o$ + FFN) & $0.922$ & $8.63$ & $-0.10$ & $+1.5\%$ \\
\bottomrule
\end{tabular}
\end{table}

Full attention-projection OV-Freeze minimises PPL (8.62) and variance
drift ($+18.2\%\rightarrow+1.3\%$); adding FFN yields no further gain.
Notably, adding OVF after inserting LayerNorm post-projection
(not shown) yields $F_1=0.921$, $\Delta\text{PPL}=-0.09$, confirming
that OV-Freeze provides complementary variance control beyond what
normalisation layers alone achieve.

\begin{table}[htbp]
\caption{OV-Freeze variance estimator ablation}
\label{tab:ovf_var_ablation}
\centering
\renewcommand{\arraystretch}{1.15}
\small
\begin{tabular}{l c c >{\centering\arraybackslash}p{1.8cm}}
\toprule
\textbf{Estimator} & \textbf{$F_1$} & \textbf{PPL amp.} & \textbf{Stability} \\
\midrule
Batch-level         & $0.918\pm0.011$ & $+1.4$ & Freq.\ oscillation \\
\textbf{EMA, $\rho=0.95$}
  & $\mathbf{0.923\pm0.006}$ & $\mathbf{+0.18}$ & \textbf{Monotone} \\
EMA, $\rho=0.90$    & $0.922\pm0.007$ & $+0.31$ & Stable \\
EMA, $\rho=0.99$    & $0.921\pm0.008$ & $+0.12$ & Slower convergence \\
Global static       & $0.919\pm0.008$ & $+0.05$ & Teacher--student drift \\
\bottomrule
\end{tabular}
\end{table}

\begin{table}[htbp]
\caption{OV-Freeze activation step-ratio ablation}
\label{tab9}
\centering
\renewcommand{\arraystretch}{1.15}
\small
\begin{tabular}{>{\raggedright\arraybackslash}p{2.5cm}ccc}
\toprule
\textbf{Activation ratio} & \textbf{$F_1$} & \textbf{PPL}
  & \textbf{Optimiser stability} \\
\midrule
0\% (no OVF)       & $0.916$ & $8.73$ & Smooth \\
10\%               & $0.919$ & $8.68$ & Smooth \\
20\%               & $0.921$ & $8.65$ & Smooth \\
\textbf{30\% (default)} & $\mathbf{0.923}$ & $\mathbf{8.62}$ & \textbf{Smooth} \\
40\%               & $0.922$ & $8.63$ & Smooth \\
50\%               & $0.918$ & $8.66$ & Severe gradient oscillation \\
\bottomrule
\end{tabular}
\end{table}

The 30\% activation window is Pareto-optimal; $\ge50\%$ causes
gradient conflict consistent with Theorem~A.2: when OV-Freeze
activates before $\mathrm{Var}_{\mathrm{EMA}}(\bm{y}_\ell)\approx
\sigma^2_{\mathrm{BF16},\ell}$, $c_\ell$ acts as a large constant
multiplier of the KL gradient, violating $\eta\le 1/\mathcal{L}'$.

\subsection{On-device Speculative Decoding}
\label{sec:speculative_exp}

\begin{table}[htbp]
\caption{Speculative decoding speedup}
\label{tab:speculative_decoding}
\centering
\renewcommand{\arraystretch}{1.2}
\small
\setlength{\tabcolsep}{3pt}
\begin{tabular}{l>{\centering\arraybackslash}p{0.45cm}>{\centering\arraybackslash}p{0.55cm}%
                >{\centering\arraybackslash}p{0.95cm}>{\centering\arraybackslash}p{0.85cm}%
                >{\centering\arraybackslash}p{0.85cm}}
\toprule
\textbf{Draft model} & $\gamma$ & $\alpha$ & \textbf{Theory} & \textbf{H100} & \textbf{SD8G3} \\
\midrule
Generic Qwen2-0.1B  & 5 & 0.78 & $3.52\times$ & $2.92\times$ & $2.78\times$ \\
\textbf{Domain-tuned (ours)} & \textbf{5} & \textbf{0.86} & $\mathbf{4.25\times}$ & $\mathbf{3.49\times}$ & $\mathbf{3.32\times}$ \\
Domain-tuned & 3 & 0.86 & $3.24\times$ & $2.65\times$ & $2.52\times$ \\
Domain-tuned & 7 & 0.86 & $5.01\times$ & $4.10\times$ & $3.90\times$ \\
Domain-tuned & 10 & 0.86 & $5.78\times$ & $4.74\times$ & $4.51\times$ \\
\bottomrule
\end{tabular}
\end{table}

Theoretical speedup ($4.25\times$) exceeds measured by 18--22\% due
to draft inference overhead ($\sim$4.8~ms/pass on SD8G3), KV-cache
management, and parallel verification dispatch ($\sim$2--3~ms/iter).
$\gamma=7$ provides only $+17\%$ throughput at $40\%$ KV-cache growth;
$\gamma=10$ pushes SD8G3 P50 to $\sim$540~ms (exceeds 500~ms budget).

\subsection{GLO Privacy Verification}
\label{sec:privacy_exp}

\begin{table}[htbp]
\caption{Speech reconstruction quality under adversarial attacks
(100-segment test set; attack models trained on 200 disjoint segments)}
\label{tab:privacy_attack}
\centering
\renewcommand{\arraystretch}{1.25}
\small
\begin{tabular}{l>{\centering\arraybackslash}p{1.5cm}>{\centering\arraybackslash}p{1.7cm}%
                >{\centering\arraybackslash}p{1.4cm}}
\toprule
\textbf{Metric} & \textbf{White-box GLO} & \textbf{Black-box U-Net} & \textbf{Random} \\
\midrule
PESQ (1--5)             & $1.21$ & $1.16$ & $1.05$ \\
WER                     & $\ge 0.95$ & $\ge 0.97$ & $1.00$ \\
MOS (1--5)              & $1.18$ & $1.11$ & $1.00$ \\
Speaker-ID accuracy     & $8.3\%$ & $7.9\%$ & $10.0\%$ \\
\bottomrule
\end{tabular}
\end{table}

Under both attack types, reconstructed-speech WER approaches the
complete-collapse upper bound ($\ge 0.95$).  Speaker-ID accuracy
(7.9--8.3\%) falls below the 10\% random baseline.  This provides
empirical evidence that the residual mutual information $I(x;
\bm{F}_v)>0$ is insufficient for any evaluated inverse attack to
recover intelligible speech.  Limitations: the attack set covers
GLO and U-Net model inversion; diffusion-based reconstruction attacks
(emerging post-training cutoff) are not evaluated and constitute a
direction for future validation.  The privacy guarantee aligns with
PIPL Article~23 data-minimisation principles ($C_1$ satisfied).

% ============================================================================
% 5. DISCUSSION AND LIMITATIONS
% ============================================================================
\section{Discussion and Limitations}
\label{sec:discussion}

\textbf{Multimodal mechanism.}  $F_1=0.923$ on TAF-28k stems from
low-bit compactness and cross-modal complementarity: acoustic
paralinguistic features (speech-rate anomalies, anxiety tones)
compensate where text-only signals are ambiguous.  OV-Freeze
maintains KL $\le 0.005$ against the BF16 teacher.

\textbf{Privacy alignment.}  Raw audio is feature-extracted on-device
and discarded; only $\bm{F}_v$ is transmitted.  GLO attack failure
(PESQ 1.21, WER $\ge 0.95$) supports practical non-reconstructability.
The $\epsilon$-LDP extension (Gaussian noise injection,
$(\epsilon,\delta)$-DP) is provided as an architectural option for
extreme-privacy scenarios; its empirical $F_1$ impact is not
evaluated in this work and is left as future work.

\textbf{Contribution boundary.}  We explicitly clarify that the
pure-KL homologous distillation paradigm is \emph{adopted} from the
NVIDIA Nemotron technical report~\cite{3}.  Our novel contributions
are (i) OV-Freeze and its formal analysis, (ii) the empirically
non-reconstructable $\bm{F}_v$ construction and acoustic proxy
statistics (Eq.~\eqref{eq:ev}--\eqref{eq:pr}), and (iii) the
end-to-end anti-fraud domain validation of this combined system.

\textbf{Misclassification analysis.}  Of 213 misclassified samples:
$\sim$41\% involve audio $<8$~s; white-box PGD-20 ($\epsilon=0.10$)
causes monotone accuracy degradation; extreme noise (SNR $<5$~dB)
induces $\sim$4.2\% discrimination drop.

\textbf{Distribution-shift assessment.}  TAF-28k is a laboratory
re-enactment distribution; results represent an empirical upper bound
under the current controlled setup.  Future work will pursue
semi-supervised domain adaptation and field deployment to quantify
the real-world generalisation gap.  NVFP4 measurements will be
replaced with native Blackwell hardware results once available.

\textbf{Latency conditions.}  Reported 268~ms P50 is under
single-call load on a lab 5G/Wi-Fi~6 network.  Under 50 simultaneous
calls, P50 degrades to $\sim$310~ms (P99 $\sim$410~ms), remaining
within the 500~ms budget.

% ============================================================================
% 6. CONCLUSION
% ============================================================================
\section{Conclusion}
\label{sec:conclusion}

This paper presents QAD-MultiGuard, an edge-side multimodal anti-fraud
framework integrating four validated pipeline components.
(1)~Pure-KL QAD (adopted from~\cite{3} and validated on TAF-28k:
$F_1=0.916$, KL 0.005).
(2)~OV-Freeze regularization (novel contribution: $+0.7$ $F_1$,
$p<0.01$), stabilising quantization-sensitive attention projections.
(3)~A 128-d empirically non-reconstructable acoustic embedding
$\bm{F}_v$ with explicitly defined proxy statistics, validated
against GLO and model-inversion attacks (WER $\ge 0.95$).
(4)~Domain-adapted speculative decoding: $3.32\times$ on SD8G3,
end-to-end P50 268~ms (single call), P50 $\sim$310~ms (50 calls).
On TAF-28k, the full pipeline achieves $F_1=0.923$, recovering
99.1\% of the BF16 upper bound.

The 248~MB quantized model achieves accuracy comparable to SAFE-QAQ
(7B, BF16) with $\sim$28$\times$ storage reduction---a
scale--efficiency trade-off.

\textbf{Limitations and future work.}  (i) Short-call robustness
($<8$~s audio) remains a gap.  (ii) Fusion weights were optimised on
TAF-28k; cross-domain recalibration should be explored.
(iii)~Diffusion-based speech reconstruction attacks should be
evaluated against $\bm{F}_v$.  (iv) Semi-supervised domain
adaptation and field deployment are the critical next steps.

% ============================================================================
% APPENDIX A — FORMAL ANALYSIS
% ============================================================================
\appendix
\section{Analysis of OV-Freeze and Acoustic Embedding Non-reconstructability}

\subsection{Proposition A.1: Boundedness of the OV-Freeze
correction factor}

\textbf{Proposition A.1.}  \emph{Let $\bm{y}_\ell\in\mathbb{R}^{N\times d}$
be the original output of layer $\ell$ and $\bm{y}'_\ell = c_\ell
\bm{y}_\ell$ the OV-Freeze-corrected output, with correction factor}
\[
  c_\ell = \sqrt{\frac{\sigma^2_{\mathrm{BF16},\ell}}
                      {\mathrm{Var}(\bm{y}_\ell)+\epsilon}},
  \quad \epsilon=10^{-6}.
\]
\emph{Then $c_\ell \in [c_{\min}, c_{\max}] \subset (0,\infty)$
for any finite batch, bounded away from zero and infinity.}

\textit{Proof.}  $\sigma^2_{\mathrm{BF16},\ell}>0$ is fixed.
$\mathrm{Var}(\bm{y}_\ell)<\infty$ for any finite batch, and
$\epsilon>0$ ensures the denominator is strictly positive.  Hence
$c_\ell \le \sigma_{\mathrm{BF16},\ell}/\sqrt{\epsilon}$ (upper
bound) and $c_\ell \ge \sigma_{\mathrm{BF16},\ell}/\sqrt{
\mathrm{Var}_{\max}+\epsilon}$ (lower bound), where
$\mathrm{Var}_{\max}$ is bounded by the BF16 teacher's variance
scale. $\square$

\textbf{Gradient boundedness under full backpropagation.}
Under the operational detached formulation~\eqref{eq:ovf-detach},
$c_\ell$ is a constant.  If the variance branch were \emph{not}
detached, the full gradient would include an additional Jacobian term:
\[
  \frac{\partial L_{\mathrm{QAD}}}{\partial\bm{y}_\ell}
    = c_\ell\!\cdot\!\frac{\partial L}{\partial\bm{y}'_\ell}
    + \left\langle\frac{\partial L}{\partial\bm{y}'_\ell},\bm{y}_\ell\right\rangle
      \!\cdot\!\frac{\partial c_\ell}{\partial\bm{y}_\ell},
\]
where
\[
  \left\|\frac{\partial c_\ell}{\partial\bm{y}_\ell}\right\|_2
    \le \frac{\sigma_{\mathrm{BF16},\ell}}{2(\mathrm{Var}+\epsilon)^{3/2}}
        \cdot\frac{2\|\bm{y}_\ell\|_2}{N-1},
\]
ruling out gradient divergence even without detachment.  The detached
implementation is a conservative simplification that inherits this
boundedness while reducing gradient complexity.

\subsection{Theorem A.2: Lipschitz convergence of the joint loss}

\textbf{Setup.}  Let $L_{\mathrm{KL}}$ have Lipschitz-continuous
gradient with constant $\mathcal{L}_0$.  OV-Freeze activates only
in the final 30\% of training, at which point EMA initialisation
ensures $c_\ell \approx 1$ for most layers.  Denote $c_{\max} =
\max_\ell\|c_\ell\|_\infty$.

\textbf{Theorem A.2.}  \emph{If $\eta \le 1/\mathcal{L}'$ with
$\mathcal{L}'=\mathcal{L}_0\cdot c_{\max}$, then SGD iterates satisfy}
\[
  \min_{t\le T}\mathbb{E}\!\left[\|\nabla L_{\mathrm{joint}}(\theta_t)\|^2\right]
    \le \frac{2(L_{\mathrm{joint}}(\theta_0)-L^\star)}{\eta T}.
\]

\textit{Proof sketch.}  OV-Freeze scales only the four
$\{q,k,v,o\}$-proj gradients by bounded constants $c_\ell$.  The
effective Lipschitz constant $\mathcal{L}' = \mathcal{L}_0 \cdot
c_{\max}$ is a conservative upper bound (the true effective constant
is smaller since $c_\ell\approx 1$ after EMA convergence and scaling
applies only to a parameter subset).  Substituting into the standard
SGD convergence bound~\cite{bubeck2015} yields the result.  This
establishes that OV-Freeze does not degrade the asymptotic convergence
rate of the underlying KL objective. $\square$

\textbf{Discussion.}  Table~\ref{tab9} confirms the theoretical
prediction: premature activation ($\ge 50\%$) causes $c_\ell$ to act
as a large multiplier before $\mathrm{Var}_{\mathrm{EMA}}\approx
\sigma^2_{\mathrm{BF16}}$, violating $\eta\mathcal{L}'\ll 1$ and
inducing gradient oscillation.

\subsection{Empirical non-reconstructability analysis of $\bm{F}_v$}

\textbf{Proposition A.3.}  \emph{Under white-box GLO and black-box
model-inversion attacks (attack models trained on 200 disjoint
segments; evaluated on 100 held-out segments), reconstructed speech
from $\bm{F}_v$ satisfies $\mathrm{WER}\ge 0.95$, $\mathrm{PESQ}
\le 1.21$, and speaker-ID accuracy $\le 8.3\%$ (random baseline:
10\%).}

\textbf{Rationale.}  Practical non-reconstructability arises from two
independent temporal compression operations: (i)~MFCC global temporal
averaging ($\sim$300 frames $\rightarrow$ one 64-d vector) discards
temporal ordering, phase, and prosodic timing; (ii)~Whisper global
average-pooling plus rank-limited projection
($\mathrm{rank}(\bm{W}_{\mathrm{proj}})\le 64$) eliminates
token-level fine-grained context.

\textbf{Important caveat.}  This is an \emph{empirical}
non-reconstructability analysis, not a formal information-theoretic
proof.  Because $\bm{F}_v$ retains task-relevant acoustic statistics
(energy variance, pitch range), $I(x;\bm{F}_v)>0$ necessarily.  The
empirical evidence indicates this residual information is insufficient
for any \emph{known} attack to recover intelligible speech or
identifiable speaker characteristics.  Diffusion-based reconstruction
attacks are not evaluated and may pose a stronger threat.  The
analysis supports PIPL Article~23 data-minimisation alignment for
practical deployment under the evaluated threat model.

% ============================================================================
% DATA AVAILABILITY
% ============================================================================
\section*{Data Availability}

\textbf{TAF-28k:} Publicly available on HuggingFace
(\texttt{JimmyMa99/TeleAntiFraud}); arXiv:2503.24115.
\textbf{ChiFraud:} Publicly available (\texttt{xuemingxxx/ChiFraud});
ACL Anthology 2025.coling-main.398.
\textbf{AdvFraud-3k:} Internal evaluation set; construction protocol
available in supplementary material.  All experiment scripts will be
released upon acceptance.

% ============================================================================
% REFERENCES
% ============================================================================
\bibliographystyle{IEEEtran}
\begin{thebibliography}{99}

\bibitem{1}
Z.~Ma, P.~Wang, M.~Huang, J.~Wang, K.~Wu, X.~Lv, Y.~Pang, Y.~Yang,
W.~Tang, and Y.~Kang, ``TeleAntiFraud-28k: An audio-text slow-thinking
dataset for telecom fraud detection,'' arXiv:2503.24115, 2025.
[Online]. Available: \url{https://arxiv.org/abs/2503.24115}

\bibitem{2}
P.~Wang \textit{et al.}, ``SAFE-QAQ: An RL-based audio-text
multimodal anti-fraud detection framework,'' arXiv preprint, 2025.
[\textit{Note: This is a preprint; production deployment figures are
as reported by the authors without independent verification.}]

\bibitem{3}
NVIDIA Research, ``NVFP4 quantization-aware distillation: Nemotron
technical report,'' NVIDIA Corp., Technical Report, Jan.\ 2026.

\bibitem{4}
G.~Hinton, O.~Vinyals, and J.~Dean, ``Distilling the knowledge in a
neural network,'' arXiv:1503.02531, 2015.

\bibitem{5}
A.~Polino, R.~Pascanu, and D.~Alistarh, ``Model compression via
distillation and quantization,'' in \emph{Proc.\ ICLR}, 2018.

\bibitem{6}
A.~Mishra and D.~Marr, ``Apprentice: Using knowledge distillation
techniques to improve low-precision network accuracy,''
arXiv:1711.05852, 2017.

\bibitem{7}
Z.~Liu \textit{et al.}, ``LLM-QAT: Data-free quantization aware
training for large language models,'' arXiv:2305.17888, 2023.

\bibitem{8}
Y.~Du \textit{et al.}, ``BitDistiller: Unleashing the potential of
sub-4-bit LLMs via self-distillation,'' in \emph{Proc.\ ACL}, 2024.

\bibitem{9}
J.~Lin \textit{et al.}, ``AWQ: Activation-aware weight quantization
for LLM compression and acceleration,'' arXiv:2306.00978, 2023.

\bibitem{10}
E.~Frantar \textit{et al.}, ``GPTQ: Accurate post-training
quantization for generative pre-trained transformers,''
arXiv:2210.17323, 2022.

\bibitem{11}
Z.~Liu \textit{et al.}, ``SpinQuant: LLM quantization with learned
rotations,'' arXiv:2405.16406, 2024.

\bibitem{12}
S.~Ashkboos \textit{et al.}, ``QuaRot: Outlier-free 4-bit inference
in rotated LLMs,'' arXiv:2404.00456, 2024.

\bibitem{13}
TODO(author): Early anti-fraud rule engine reference 1.

\bibitem{14}
TODO(author): Early anti-fraud rule engine reference 2.

\bibitem{15}
TODO(author): Early anti-fraud rule engine reference 3.

\bibitem{16}
TODO(author): BERT-Fraud reference ($F_1 = 0.876$, Chinese SMS dataset).

\bibitem{17}
G.~Gerganov \textit{et al.}, ``GGUF: A format for storing and
distributing large language models,'' llama.cpp project, 2023.
[Online]. Available: \url{https://github.com/ggerganov/llama.cpp}

\bibitem{18}
W.~Kwon \textit{et al.}, ``Efficient memory management for large
language model serving with PagedAttention,'' in \emph{Proc.\ SOSP},
2023.

\bibitem{19}
W.~Kwon \textit{et al.}, ibid.\ (vLLM, PagedAttention---same
publication as~\cite{18}).

\bibitem{20}
Y.~Leviathan, M.~Kalman, and Y.~Matias, ``Fast inference from
transformers via speculative decoding,'' in \emph{Proc.\ ICML}, 2023.

\bibitem{21}
C.~Chen \textit{et al.}, ``Accelerating large language model decoding
with speculative sampling,'' arXiv:2302.01318, 2023.

\bibitem{22}
N.~Tomashenko \textit{et al.}, ``Introducing the VoicePrivacy
initiative,'' in \emph{Proc.\ Interspeech}, 2020.

\bibitem{23}
A.~Bora \textit{et al.}, ``Compressed sensing using generative
models,'' in \emph{Proc.\ ICML}, 2017. [\textit{Used in this work to
motivate the GLO attack threat model.}]

\bibitem{24}
TODO(author): QAT reference 1.

\bibitem{25}
TODO(author): QAT reference 2.

\bibitem{26}
TODO(author): QAT reference 3.

\bibitem{27}
J.~Yang \textit{et al.}, ``Qwen2.5 technical report,''
arXiv:2412.15115, 2024.

\bibitem{28}
Qwen Team, ``Qwen2.5: A party of foundation models,''
arXiv:2412.15115, 2024.

\bibitem{29}
Qwen Team, ``Qwen2 technical report,'' arXiv:2407.10671, 2024.

\bibitem{bubeck2015}
S.~Bubeck, ``Convex optimization: Algorithms and complexity,''
\emph{Found.\ Trends Mach.\ Learn.}, vol.~8, no.~3--4, pp.~231--357,
2015.

\bibitem{chifraud}
M.~Tang, L.~Zou, S.-N.~Liang, Z.~Jin, W.~Wang, and S.~Cui,
``ChiFraud: A long-term web text benchmark for Chinese fraud
detection,'' in \emph{Proc.\ COLING}, 2025, pp.~5962--5974.
[Online]. Available: \url{https://aclanthology.org/2025.coling-main.398/}

\end{thebibliography}

\end{document}
